Linear stability lobes can overestimate the usable depth of cut when an active damping actuator saturates. This computational study asks how much disturbance a linearly stable milling cut with voltage-limited velocity feedback tolerates before saturation triggers chatter. In a single-mode model of a published case study (17 of 19 tests classified as measured), the command splits into a forced part and the residual authority left for disturbances, which governs the basin near the limit. An abrupt entry chatters below 97\,\% of the linear limit at 157 of 251 speeds; near the limit, kicks of 2.4--10.8\,\% of one tooth's cutting impulse end in bounded clipped chatter coexisting with the stable cut. A circle criterion, its frequency condition proved by interval arithmetic, shows the forced orbit of the averaged delay model to attract every solution below a depth whose median is 0.97 of the uncontrolled limit; above it, linear matrix inequalities give inner estimates of the basin, numerically verified for collocation models and never violated in simulation. A basin-aware gain schedule, searched on grids at 21 speeds, gives the largest depth that is linearly stable, keeps a voltage margin and is certified against a required disturbance. With equal gain range and margins its gains certify 6\,mm more depth than authority-based ones for 5\,mm/s, hardly more for 1\,mm/s; with abrupt entry or a 16\,mm cap it cuts no deeper. Under sampled model errors its points, nearer their linear limit, fail more often (10--17\,\% against 7\,\%).
